% foundations/operators.tex

\chapter{Typed Operators}
\label{chap:operators}

The revised curvature has produced a revised intrinsic state. The next
question is not which concepts exist, but which action produces each
subsequent consequence. AIM represents such actions as typed operators.

\begin{definition}[Operator]
\label{def:operator}
An operator is a mapping
\[
\mathcal O_c:\mathcal A\rightarrow\mathcal B
\]
whose input type, output type, applicable context \(c\), assumptions,
and result provenance are declared.
\end{definition}

The subscript is shown when context would otherwise be ambiguous. An
algorithm may implement an operator, but the engineering contract is not
defined by implementation alone.

\begin{principle}[State--operator relation]
\label{princ:state-operator-relation}
Operators evaluate, transform, or relate engineering states without
becoming constituents of the engineering object itself.
\end{principle}

\begin{proposition}
\label{prop:operator-explicit-dependencies}
Engineering dependencies that change the meaning, domain, or status of a
state shall be represented explicitly through operators or relations
rather than hidden inside an undifferentiated state description.
\end{proposition}

\section{Reconstruct the edited state}

The geometric reconstruction operator can be written schematically as
\[
\mathcal D_{I,s_0}:
(\mathrm{pose2}(s_0),\kappa)
\longrightarrow \mathrm{pose2}|_I.
\]

\paragraph{Execution contract.}
Input: an initial pose, curvature function, and interval. Output: the
derived state trajectory. Required context: longitudinal domain and
boundary conditions. Applicability: the reconstruction axioms and
regularity assumptions must hold. Preserved identity: the alignment to
which the edited state is assigned. Provenance: the edit, initial state,
and reconstruction method. It cannot decide admissibility or whether
the edited state should replace an approved state.

\section{Realize the state metrically}

\begin{definition}[Metric realization operator]
\label{def:metric-realization-operator}
\[
\mathcal G_c:\mathcal X\rightarrow\mathcal X_{\mathrm{metric}}.
\]
\end{definition}

\paragraph{Execution contract.}
Input: the intrinsic state. Output: its metric interpretation. Required
context: metric model, embedding data, and applicable conventions.
Applicability: the context must cover the state domain and requested
quantities. Preserved identity: the engineering object and state role.
Provenance: intrinsic source plus realization context. It cannot perform
physical construction, create an observation, or select a decision.

\begin{definition}[Frame operator]
\label{def:frame-operator}
The frame operator constructs the local frame required by the applicable
realization:
\[
\mathcal F_c:\mathcal X_{\mathrm{metric}}\rightarrow\mathcal F.
\]
\end{definition}

Its output is derived orientation data. It preserves identity and
inherits the state and realization provenance; it cannot determine which
frame convention a project ought to adopt.

\section{Relate track and world descriptions}

\begin{definition}[Track-to-world operator]
\label{def:track-to-world-operator}
\[
\mathcal P_{\mathrm{tw},c}:
\Omega_{\mathrm{track}}\rightarrow\Omega_{\mathrm{world}}.
\]
\end{definition}

It receives track coordinates and a realized alignment frame and
produces a world-space point. Its applicability requires a valid local
frame over the requested domain. The result preserves the referenced
object and carries the realization provenance; it cannot establish a
physical or observed condition.

\begin{definition}[World-to-track operator]
\label{def:world-to-track-operator}
\[
\mathcal P_{\mathrm{wt},c}:
\Omega_{\mathrm{world}}\rightarrow\Omega_{\mathrm{track}}.
\]
\end{definition}

It receives a world-space point and applicable realized alignment and
produces a candidate track-related interpretation. Existence and
uniqueness depend on the search domain and geometry. Its provenance must
retain the world source, realization, and selection rule. It cannot prove
that a survey point observes the alignment or choose among ambiguous
projections without an applicable rule.

\begin{definition}[Coordinate transformation operator]
\label{def:coordinate-transform-operator}
\[
\mathcal C_c:
\Omega_{\mathrm{world}}\rightarrow\Omega_{\mathrm{world}}.
\]
\end{definition}

It changes coordinate representation under a declared transformation
context. It preserves the represented object's identity and retains
source/target CRS provenance. It cannot replace world-to-track
interpretation or metric realization.

\section{Produce representations and observations}

Sampling the revised metric trajectory produces a CAD curve or polyline
representation. The representation operator receives qualified state
knowledge and representation parameters; it produces a derived artifact,
preserves object identity, and retains source-state and method provenance.
It cannot become the constructive definition merely because it is
stored, exchanged, or rendered.

An observation operation instead receives a physical condition through
a measurement process and produces source-grounded evidence with time,
method, conditions, and uncertainty. A survey polyline is usable in the
comparison only through that provenance. Observation cannot turn its
output into physical reality or admitted engineering knowledge by
itself.

\section{Relate target and physical condition}

\begin{definition}[Realization operator]
\label{def:realization-operator}
\[
\mathcal R:
\mathcal X_{\mathrm{target}}\rightarrow\mathcal X_{\mathrm{real}}.
\]
\end{definition}

The operator receives a target state plus an applicable construction,
maintenance, or physical-response model and produces a realized-state
description. It preserves the identity of the engineering object while
changing the condition described. Its provenance includes the target,
process model, and relevant physical inputs. It cannot manufacture an
observation of the result or decide that the realized condition is
acceptable.

\section{Evaluate the qualified consequence}

\begin{definition}[Evaluation operator]
\label{def:evaluation-operator}
\[
\mathcal E_c:
\mathcal X\times\Omega_{\mathrm{track}}\rightarrow\mathcal Y.
\]
\end{definition}

For the recurring discrepancy, evaluation receives the metrically
realized target state, compatible survey evidence, comparison locations,
uncertainty, and applicable rules. It produces residuals or other
qualified consequences. It preserves the identity and roles of its
inputs, and its provenance records those inputs, context, operator
version, and assumptions. It cannot admit evidence as knowledge, accept
a proposal, or issue an engineering decision.

\begin{principle}[Solver independence]
\label{princ:operators-solver-independence}
Engineering observations, constraints, residuals, candidate solutions,
and decisions are independent of the numerical solver used to evaluate
them.
\end{principle}

\section{Composition without responsibility transfer}

For compatible types,
\[
\mathcal O_1:\mathcal A\rightarrow\mathcal B,
\qquad
\mathcal O_2:\mathcal B\rightarrow\mathcal C,
\qquad
\mathcal O_2\circ\mathcal O_1:
\mathcal A\rightarrow\mathcal C.
\]

\begin{principle}[Operator separation]
\label{princ:operator-separation}
Operators with different engineering responsibilities shall remain
conceptually distinguishable even when they participate in the same
workflow.
\end{principle}

The curvature-edit chain can therefore be read as reconstruction,
metric realization, representation or observation interpretation, and
evaluation. Composition transmits values and provenance. It does not
transfer identity, change truth mode silently, or confer decision
authority on the final number.
